% ============================================================
%  Chapter 11 — Flexible Nuisance Estimation, Orthogonal Scores,
%               and Cross-Fitting
%  Compile standalone:  pdflatex chapter_11.tex
%  Compile via master:  pdflatex master.tex
% ============================================================
\documentclass[master.tex]{subfiles}

\begin{document}

% ------------------------------------------------------------
%  Title block (standalone) / chapter heading (master)
% ------------------------------------------------------------
\ifSubfilesClassLoaded{%
  \begin{center}
    {\LARGE\bfseries\color{isublue}
     Chapter 11: Flexible Nuisance Estimation, Orthogonal Scores,
     and Cross-Fitting}\\[6pt]
    {\large Theory and Methods in Causal Inference}\\[4pt]
    {\normalsize Department of Statistics, Iowa State University}
  \end{center}
  \bigskip
  \hrule height 1.2pt \relax
  \smallskip
}{%
  \chapter{Flexible Nuisance Estimation, Orthogonal Scores, and Cross-Fitting}
  \label{w11:chap}
}

% ------------------------------------------------------------
%  Learning objectives
% ------------------------------------------------------------
\begin{tcolorbox}[defstyle, title={Learning Objectives}]
By the end of this lecture, students should be able to:
\begin{enumerate}
  \item Explain why flexible nuisance estimation can produce misleading
    inference for a causal parameter even when predictive accuracy
    is high.
  \item Define Neyman orthogonality, verify it for a given estimating
    function, and explain how it reduces sensitivity to
    first-order nuisance estimation error.
  \item Identify the orthogonal score for the ATE, recognize it as
    the efficient influence function from Chapter~10, and state the
    three roles it simultaneously fulfills.
  \item Describe sample splitting, explain why independence between
    nuisance estimation and score evaluation simplifies
    asymptotic arguments, and articulate its efficiency cost.
  \item Describe $K$-fold cross-fitting, write the cross-fitted
    moment equation explicitly, and explain how cross-fitting
    recovers efficiency relative to a single split.
  \item State the product-rate condition for the cross-fitted AIPW
    estimator, connect it to Neyman orthogonality, and explain
    why it allows each nuisance estimator to converge at rate
    $n^{-1/4}$.
  \item Construct a consistent variance estimator from
    out-of-fold estimated influence values and form a
    Wald confidence interval.
  \item Articulate what machine learning can and cannot contribute
    to causal inference, and follow the practical workflow for
    cross-fitted orthogonal-score estimation.
\end{enumerate}
\end{tcolorbox}


% ============================================================
\section{Why Flexible Nuisance Estimation Is Both Attractive and Dangerous}
\label{w11:sec:motivation}
% ============================================================

In Chapters~9 and~10 we developed estimation and inference using
estimating equations, influence functions, doubly robust scores, and
semiparametric efficiency.  A central feature of these methods is
that the causal parameter of interest depends on nuisance functions
such as the outcome regressions
\[
  \mu_t(x) = \E(Y \mid T=t,\, X=x), \qquad t \in \{0,1\},
\]
and the propensity score
\[
  \pi(x) = P(T=1 \mid X=x).
\]
In many real applications, these nuisance functions may be too
complex to model accurately using low-dimensional parametric forms.
This makes flexible methods attractive, including penalized
regression, splines and sieves, random forests, boosting, neural
networks, and ensemble learning.

These methods can substantially improve predictive accuracy, but they
also create new statistical difficulties:
\begin{itemize}
  \item they may converge more slowly than parametric estimators;
  \item they may overfit the observed sample;
  \item their asymptotic behavior may be difficult to characterize;
  \item naive plug-in inference can fail even when prediction quality
    is good.
\end{itemize}

Better nuisance prediction therefore does not automatically imply
valid inference for the target causal parameter.  The goal of this
chapter is to explain how orthogonal scores, sample splitting, and
cross-fitting make it possible to combine flexible nuisance
estimation with valid large-sample inference.  The solution
developed here combines an orthogonal score, which removes
first-order sensitivity to nuisance estimation error, with
cross-fitting, which separates nuisance estimation from score
evaluation to control the remaining empirical process remainder.


% ============================================================
\section{Why Naive Plug-In Estimation Can Fail}
\label{w11:sec:plugin}
% ============================================================

Suppose the parameter of interest is written as a functional
$\psi = \Psi(\eta)$, where $\eta$ denotes a nuisance object such
as $(\mu_0, \mu_1, \pi)$.  A natural estimator is the
\emph{plug-in estimator}
\[
  \hat\psi_{\mathrm{plug}} = \Psi(\hat\eta),
\]
obtained by replacing the nuisance functions with estimators.

\subsection{The first-order Taylor expansion}
\label{w11:subsec:taylor}

To see precisely how nuisance error propagates, apply a
first-order Taylor expansion of $\Psi$ around the true nuisance
$\eta_0$:
\begin{align}
  \hat\psi_{\mathrm{plug}} - \psi
  =
  \underbrace{%
    D_\eta\Psi(\eta_0)[\hat\eta - \eta_0]
  }_{\text{linear term}}
  +
  \underbrace{%
    R(\hat\eta,\,\eta_0)
  }_{\text{second-order remainder}},
  \label{w11:eq:taylor}
\end{align}
where $D_\eta\Psi(\eta_0)[h]$ denotes the Gateaux derivative of
$\Psi$ at $\eta_0$ in direction $h$, and the remainder satisfies
$|R(\hat\eta, \eta_0)| \lesssim \|\hat\eta - \eta_0\|^2$.

For root-$n$ inference we need $\sqrt{n}(\hat\psi_{\mathrm{plug}}
- \psi)$ to have a centered limiting distribution.  The
second-order remainder is typically manageable: if
$\|\hat\eta - \eta_0\| = o_p(n^{-1/4})$ then
$R = o_p(n^{-1/2})$, which is negligible at the root-$n$ scale.
The obstacle is the \emph{linear term}, which is proportional to
the nuisance estimation error $\hat\eta - \eta_0$.

\subsection{Finite-dimensional nuisance: orthogonality is sufficient}
\label{w11:subsec:finite-dim}

Suppose first that $\eta \in \mathbb{R}^d$ for a fixed finite $d$,
and that $\hat\eta$ is a standard parametric estimator satisfying
$\hat\eta - \eta_0 = O_p(n^{-1/2})$.  Then the linear term in
\eqref{w11:eq:taylor} is also $O_p(n^{-1/2})$, which contaminates
the root-$n$ limiting distribution of $\hat\psi_{\mathrm{plug}}$
with a bias of order $n^{-1/2}$.

If the functional $\Psi$ satisfies $D_\eta\Psi(\eta_0) = 0$ —
the \emph{Neyman orthogonality} condition introduced formally in
Section~\ref{w11:sec:orthogonal} — then the linear term vanishes
identically.  The expansion \eqref{w11:eq:taylor} reduces to the
second-order remainder alone, which under parametric rates is
$O_p(n^{-1})$ and hence asymptotically negligible.  In the
finite-dimensional setting, \textbf{Neyman orthogonality is
sufficient} to eliminate the leading source of bias.

\subsection{Infinite-dimensional nuisance: orthogonality is not enough}
\label{w11:subsec:infinite-dim}

When $\eta$ belongs to a function space — as is the case for the
nuisance triple $(\mu_0, \mu_1, \pi)$ — orthogonality is no longer
sufficient on its own.  Orthogonality removes the first derivative
of the target with respect to nuisance perturbations, but it does
not by itself control the stochastic fluctuation of the score when
the nuisance estimate is a random function chosen from a large
class.  A second source of first-order error arises from evaluating
the estimating equation on the same sample used to estimate
$\hat\eta$.  Write the
plug-in estimating equation as
$\mathbb{P}_n\{\varphi(\,\cdot\,;\psi, \hat\eta)\} = 0$, and
decompose the deviation from the population equation as
\begin{align}
  &\mathbb{P}_n\{\varphi(\,\cdot\,;\psi, \hat\eta)\}
  -
  \mathbb{P}_n\{\varphi(\,\cdot\,;\psi, \eta_0)\}
  \nonumber\\[4pt]
  &\quad=
  \underbrace{%
    P\bigl\{\varphi(\,\cdot\,;\psi,\hat\eta)
       - \varphi(\,\cdot\,;\psi,\eta_0)\bigr\}
  }_{\text{population bias}}
  +
  \underbrace{%
    (\mathbb{P}_n - P)\bigl\{\varphi(\,\cdot\,;\psi,\hat\eta)
                        - \varphi(\,\cdot\,;\psi,\eta_0)\bigr\}
  }_{\text{empirical process term}}.
  \label{w11:eq:decomp}
\end{align}
Neyman orthogonality controls the \emph{population bias}: because
the Gateaux derivative of $P\varphi$ with respect to $\eta$
vanishes at the truth, the population bias is second order in
$\|\hat\eta - \eta_0\|$.

However, Neyman orthogonality says nothing about the
\emph{empirical process term}.  This term depends on the
complexity of the function class
$\{\varphi(\,\cdot\,;\psi,\eta) : \eta \in \mathcal{H}\}$
traced out as $\hat\eta$ varies.  For parametric estimators the
class is small, and standard empirical process theory controls
the term at order $o_p(n^{-1/2})$.  For flexible machine-learning
estimators — random forests, neural networks, high-dimensional
lasso — the class may be too large for such arguments to apply
(the \emph{Donsker condition} fails), and the empirical process
term can remain non-negligible even as $\hat\eta$ converges
consistently.

The essential conclusion is:
\begin{quote}
  In the infinite-dimensional setting, Neyman orthogonality
  eliminates the population-level linear bias but does not
  control the empirical process remainder.  Valid root-$n$
  inference requires both: orthogonality to remove the
  population bias, and \emph{sample splitting} or
  \emph{cross-fitting} to remove the empirical process term.
\end{quote}
The remainder of the chapter addresses these two problems in turn:
\begin{itemize}
  \item \textbf{Problem~1 (population bias).}  Controlled by
    Neyman orthogonality.  Section~\ref{w11:sec:orthogonal}
    defines orthogonality formally and
    Section~\ref{w11:sec:ate-score} verifies it for the ATE score.
  \item \textbf{Problem~2 (empirical process term).}  Not
    controlled by orthogonality alone.  Section~\ref{w11:sec:reuse}
    makes this term precise and explains when it fails to vanish;
    Sections~\ref{w11:sec:splitting} and~\ref{w11:sec:crossfit}
    show how sample splitting and cross-fitting eliminate it.
\end{itemize}

\begin{remark}
  This problem is not specific to machine learning.  It arises
  whenever nuisance parameters are estimated in an
  infinite-dimensional model.  Machine learning makes the issue
  more visible because it encourages highly adaptive nuisance
  estimation; see \citet{chernozhukov2018double} for a systematic
  treatment.
  % BIBLIOGRAPHY NOTE: chernozhukov2018double is not yet in
  % causal_inference.bib.  Add manually:
  %   Chernozhukov, V., Chetverikov, D., Demirer, M., Duflo, E.,
  %   Hansen, C., Newey, W., and Robins, J. (2018).
  %   Double/debiased machine learning for treatment and structural
  %   parameters.  The Econometrics Journal, 21(1), C1--C68.
\end{remark}


% ============================================================
\section{Orthogonal Scores}
\label{w11:sec:orthogonal}
% ============================================================

The key device for reducing sensitivity to nuisance estimation error
is an \emph{orthogonal score}.  Let
\[
  \E\{\phi(O;\,\psi,\eta)\} = 0
\]
be an estimating equation for the target parameter $\psi$, where
$\eta$ denotes nuisance parameters.

\begin{defn}{Neyman Orthogonality}
\label{w11:defn:orthogonal}
The score function $\phi(O;\,\psi,\eta)$ is called \emph{orthogonal}
(or \emph{Neyman orthogonal}) at the truth $(\psi_0, \eta_0)$ if the
Gateaux derivative of the estimating equation with respect to the
nuisance, evaluated at the truth, vanishes:
\[
  \left.
    \frac{\partial}{\partial r}
    \E\bigl\{\phi(O;\,\psi_0,\,\eta_0 + r h)\bigr\}
  \right|_{r=0}
  = 0
\]
for all perturbation directions $h$ in a suitable class.
\end{defn}

This condition means that small local perturbations of the nuisance
parameter have no first-order effect on the estimating equation at the
truth.  Consequently, nuisance estimation error enters only at second
order, at least locally.  Orthogonality does not mean nuisance
estimation no longer matters; it means nuisance error enters only
through a second-order population remainder, provided the score is
evaluated in a way that does not introduce additional same-sample
overfitting bias.

Orthogonality is the principal reason that semiparametric estimators
can remain valid under flexible nuisance estimation.  In particular,
the doubly robust score introduced in Chapter~10 is orthogonal, which
is exactly why the product-rate condition in
\eqref{w10:eq:product-rate} is sufficient for asymptotic negligibility
of nuisance estimation error.

\begin{remark}[Neyman's original insight]
\label{w11:rem:neyman-history}
The terminology honors Jerzy Neyman, whose work on hypothesis testing
introduced the idea of constructing test statistics that are insensitive
to nuisance parameters.  In the modern semiparametric literature the
concept was formalized by \citet{robins1994estimation} and others,
and later made systematic in the double machine learning framework
of \citet{chernozhukov2018double}.
\end{remark}


% ============================================================
\section{The Orthogonal Score for the Average Treatment Effect}
\label{w11:sec:ate-score}
% ============================================================

We continue to use the average treatment effect
\[
  \tau = \E\{Y(1) - Y(0)\}
\]
as the running example.  Under consistency, conditional
exchangeability, and positivity, the efficient influence function
derived in Chapter~10 (equation~\eqref{w10:eq:eif-ate}) is
\begin{align}
  \varphi_{\mathrm{eff}}(O;\,\tau,\eta)
  =
  \frac{T}{\pi(X)}\{Y - \mu_1(X)\}
  -
  \frac{1-T}{1-\pi(X)}\{Y - \mu_0(X)\}
  +
  \mu_1(X) - \mu_0(X) - \tau,
  \label{w11:eq:orth-score-ate}
\end{align}
where $\eta = (\mu_0, \mu_1, \pi)$.  Chapter~10 identified this as
the right score for the ATE; the role of Chapter~11 is to explain
how to estimate its nuisance components safely when flexible learners
are used.  DML is not a new estimator but a new nuisance-fitting
protocol built around the same score.  We now prove that this score
is Neyman orthogonal.  The proof verifies
Definition~\ref{w11:defn:orthogonal} separately for perturbations
in each of the three nuisance components.

\begin{lemma}{Neyman Orthogonality of the ATE Score}
\label{w11:lem:orthogonal-ate}
The score $\varphi_{\mathrm{eff}}(O;\,\tau,\eta)$ in
\eqref{w11:eq:orth-score-ate} is Neyman orthogonal at every truth
$(\tau_0, \eta_0) = (\tau_0, \mu_0, \mu_1, \pi)$ satisfying the
identification assumptions.  That is, for every bounded perturbation
direction $h(x)$, $g(x)$, or $\ell(x)$,
\begin{align*}
  \left.\frac{\partial}{\partial r}
    \E\!\left\{\varphi_{\mathrm{eff}}(O;\,\tau_0,\,\mu_0,\,\mu_1,\,\pi+rh)\right\}
  \right|_{r=0} &= 0, \\[4pt]
  \left.\frac{\partial}{\partial r}
    \E\!\left\{\varphi_{\mathrm{eff}}(O;\,\tau_0,\,\mu_0,\,\mu_1+rg,\,\pi)\right\}
  \right|_{r=0} &= 0, \\[4pt]
  \left.\frac{\partial}{\partial r}
    \E\!\left\{\varphi_{\mathrm{eff}}(O;\,\tau_0,\,\mu_0+r\ell,\,\mu_1,\,\pi)\right\}
  \right|_{r=0} &= 0.
\end{align*}
\end{lemma}

\begin{proof}
Throughout, differentiation and expectation are interchanged under
standard bounded-convergence arguments, and the law of iterated
expectations is applied by conditioning on $X$.  Recall that under
conditional exchangeability and consistency,
\[
  \E(Y \mid T=1, X) = \mu_1(X), \qquad
  \E(Y \mid T=0, X) = \mu_0(X), \qquad
  \E(T \mid X) = \pi(X).
\]

\medskip
\noindent\textit{Perturbation in $\pi$.}\enspace
Replace $\pi$ by $\pi + rh$ and differentiate.  The only terms in
\eqref{w11:eq:orth-score-ate} that depend on $\pi$ are the two
IPW residual terms, so
\begin{align*}
  \frac{\partial}{\partial r}
  \E\!\left\{\varphi_{\mathrm{eff}}(O;\,\tau_0,\mu_0,\mu_1,\pi+rh)\right\}
  \bigg|_{r=0}
  &=
  \E\!\left[
    -\frac{T\,h(X)}{\pi(X)^2}\{Y-\mu_1(X)\}
    -\frac{(1-T)\,h(X)}{(1-\pi(X))^2}\{Y-\mu_0(X)\}
  \right].
\end{align*}
Conditioning on $X$ and using $\E(T \mid X) = \pi(X)$,
\begin{align*}
  \E\!\left[T\{Y-\mu_1(X)\}\mid X\right]
  &= \E(TY \mid X) - \mu_1(X)\,\pi(X) \\
  &= \pi(X)\,\mu_1(X) - \mu_1(X)\,\pi(X) = 0, \\[4pt]
  \E\!\left[(1-T)\{Y-\mu_0(X)\}\mid X\right]
  &= \E((1-T)Y \mid X) - \mu_0(X)(1-\pi(X)) \\
  &= (1-\pi(X))\,\mu_0(X) - \mu_0(X)(1-\pi(X)) = 0.
\end{align*}
Both conditional expectations vanish, so the derivative is zero.

\medskip
\noindent\textit{Perturbation in $\mu_1$.}\enspace
Replace $\mu_1$ by $\mu_1 + rg$.  The terms in
\eqref{w11:eq:orth-score-ate} that depend on $\mu_1$ are
$\frac{T}{\pi(X)}\{Y-\mu_1(X)\}$ and $\mu_1(X)$.  Differentiating,
\begin{align*}
  \frac{\partial}{\partial r}
  \E\!\left\{\varphi_{\mathrm{eff}}(O;\,\tau_0,\mu_0,\mu_1+rg,\pi)\right\}
  \bigg|_{r=0}
  &=
  \E\!\left[-\frac{T}{\pi(X)}\,g(X) + g(X)\right] \\
  &=
  \E\!\left[g(X)\left(1 - \frac{T}{\pi(X)}\right)\right].
\end{align*}
Conditioning on $X$,
\[
  \E\!\left[1 - \frac{T}{\pi(X)}\;\middle|\; X\right]
  = 1 - \frac{\E(T\mid X)}{\pi(X)}
  = 1 - \frac{\pi(X)}{\pi(X)} = 0,
\]
so the derivative is zero.

\medskip
\noindent\textit{Perturbation in $\mu_0$.}\enspace
Replace $\mu_0$ by $\mu_0 + r\ell$.  The terms in
\eqref{w11:eq:orth-score-ate} that depend on $\mu_0$ are
$-\frac{1-T}{1-\pi(X)}\{Y-\mu_0(X)\}$ and $-\mu_0(X)$.
Differentiating,
\begin{align*}
  \frac{\partial}{\partial r}
  \E\!\left\{\varphi_{\mathrm{eff}}(O;\,\tau_0,\mu_0+r\ell,\mu_1,\pi)\right\}
  \bigg|_{r=0}
  &=
  \E\!\left[\frac{1-T}{1-\pi(X)}\,\ell(X) - \ell(X)\right] \\
  &=
  \E\!\left[\ell(X)\left(\frac{1-T}{1-\pi(X)} - 1\right)\right].
\end{align*}
Conditioning on $X$,
\[
  \E\!\left[\frac{1-T}{1-\pi(X)} - 1\;\middle|\; X\right]
  = \frac{1 - \E(T\mid X)}{1-\pi(X)} - 1
  = \frac{1-\pi(X)}{1-\pi(X)} - 1 = 0,
\]
so the derivative is zero.  This completes the proof.
\end{proof}

\begin{remark}[What the proof reveals]
\label{w11:rem:orthog-insight}
Each of the three derivative calculations reduces to a single key
fact: the conditional mean of the IPW weight equals one.
Specifically, $\E(T/\pi(X)\mid X) = 1$ and
$\E((1-T)/(1-\pi(X))\mid X) = 1$ are exactly the balancing
properties of the propensity score established in
Chapter~6 (Theorem on the Balancing Property).  Neyman
orthogonality of the ATE score is therefore a direct consequence
of propensity score balancing.
\end{remark}

The score in \eqref{w11:eq:orth-score-ate} simultaneously fulfills
three purposes:
\begin{enumerate}
  \item it identifies $\tau$ through the moment condition
    $\E\{\varphi_{\mathrm{eff}}(O;\,\tau,\eta)\} = 0$;
  \item it is doubly robust, yielding a consistent estimator
    whenever either the outcome model or the propensity score model
    is correctly specified; and
  \item it is the efficient influence function, so any estimator
    based on it achieves the semiparametric efficiency bound
    (Theorem~\ref{w10:thm:bound}).
\end{enumerate}
This is why the score in \eqref{w11:eq:orth-score-ate} is central
in modern causal inference.


% ============================================================
\section{Why Reusing the Same Data Can Be Problematic}
\label{w11:sec:reuse}
% ============================================================

Section~\ref{w11:sec:plugin} identified two obstacles to valid
root-$n$ inference: the population bias and the empirical process
term in decomposition~\eqref{w11:eq:decomp}.  We showed there that
Neyman orthogonality eliminates the population bias but leaves the
empirical process term uncontrolled.  This section makes the
empirical process term precise in the context of the ATE estimating
equation and introduces the Donsker condition — the classical tool
for controlling it when the same data are reused for both nuisance
estimation and score evaluation.

\subsection{The plug-in remainder for the AIPW estimator}
\label{w11:subsec:ep-remainder}

Recall from Chapter~9 that we write $\mathbb{P}_n f =
n^{-1}\sum_{i=1}^n f(O_i)$ for the empirical average of a function
$f$, and define the \emph{centered empirical process}
\[
  \mathbb{G}_n f
  = \sqrt{n}\,(\mathbb{P}_n - P) f
  = \frac{1}{\sqrt{n}}\sum_{i=1}^n \{f(O_i) - \E f(O)\}.
\]
Consider the plug-in AIPW estimating equation
$\mathbb{P}_n\{\varphi_{\mathrm{eff}}(O;\,\tau,\hat\eta)\} = 0$,
where the same sample is used for both nuisance estimation and
score evaluation.  Expanding around $\eta_0$ gives
\begin{align}
  0
  = \mathbb{P}_n\{\varphi_{\mathrm{eff}}(O;\,\tau,\eta_0)\}
    + \underbrace{%
        \mathbb{P}_n\{
          \varphi_{\mathrm{eff}}(O;\,\tau,\hat\eta)
          -
          \varphi_{\mathrm{eff}}(O;\,\tau,\eta_0)
        \}
      }_{=:\,R_n}.
  \label{w11:eq:plugin-expand}
\end{align}
The first term is a sample average of a fixed function, satisfies
the CLT, and is $O_p(n^{-1/2})$.  The remainder $R_n$ is exactly
the left-hand side of decomposition~\eqref{w11:eq:decomp}
specialized to $\varphi = \varphi_{\mathrm{eff}}$:
\begin{align}
  R_n
  = \underbrace{%
      P\{
        \varphi_{\mathrm{eff}}(\,\cdot\,;\tau,\hat\eta)
        - \varphi_{\mathrm{eff}}(\,\cdot\,;\tau,\eta_0)
      \}
    }_{\text{population bias}}
    +
    \underbrace{%
      (\mathbb{P}_n - P)\{
        \varphi_{\mathrm{eff}}(\,\cdot\,;\tau,\hat\eta)
        - \varphi_{\mathrm{eff}}(\,\cdot\,;\tau,\eta_0)
      \}
    }_{\text{empirical process term}}.
  \label{w11:eq:Rn-decomp}
\end{align}
By Lemma~\ref{w11:lem:orthogonal-ate}, Neyman orthogonality ensures
the population bias is $o_p(n^{-1/2})$ under the product-rate
condition.  Writing $f_\eta(\cdot) =
\varphi_{\mathrm{eff}}(\,\cdot\,;\tau,\eta)$, the empirical process
term takes the form
\begin{align}
  (\mathbb{P}_n - P)\{f_{\hat\eta} - f_{\eta_0}\}
  =
  \frac{1}{\sqrt{n}}\,\mathbb{G}_n\{f_{\hat\eta} - f_{\eta_0}\}.
  \label{w11:eq:ep-term}
\end{align}
This term depends on $\hat\eta$, which is estimated from the same
sample.  Whether it is $o_p(n^{-1/2})$ depends on how complex the
class $\{f_\eta : \eta \in \mathcal{H}\}$ is — a question the
Donsker condition answers.

\subsection{The Donsker condition}
\label{w11:subsec:donsker}

The Donsker condition is a technical way of asking whether the
function class used by the nuisance learner is tame enough that
same-sample score evaluation remains asymptotically well behaved.

\begin{defn}{Donsker Class}
\label{w11:defn:donsker}
A class of measurable functions $\mathcal{F}$ is a
\emph{Donsker class} (with respect to $P$) if the empirical process
$\{\mathbb{G}_n f : f \in \mathcal{F}\}$ converges weakly in
$\ell^\infty(\mathcal{F})$ to a tight Gaussian process.
Equivalently, $\mathcal{F}$ is Donsker if
\[
  \sup_{f \in \mathcal{F}}
  \bigl|(\mathbb{P}_n - P) f\bigr|
  = O_p(n^{-1/2}),
\]
and the fluctuations of $n^{1/2}(\mathbb{P}_n - P)$ over
$\mathcal{F}$ are stochastically equicontinuous; see
\citet{vandervaart1998}, Chapter~19.
\end{defn}

A key sufficient condition for $\mathcal{F}$ to be Donsker is that
its \emph{bracketing entropy} is finite.  The bracketing number
$N_{[\,]}(\epsilon, \mathcal{F}, L_2(P))$ counts the minimum number
of $\epsilon$-brackets $[\ell_j, u_j]$ needed to cover
$\mathcal{F}$, where a bracket is a pair of functions with
$P(u_j - \ell_j)^2 \leq \epsilon^2$.  The class $\mathcal{F}$ is
Donsker whenever
\[
  \int_0^\infty
  \sqrt{\log N_{[\,]}(\epsilon, \mathcal{F}, L_2(P))}\,
  d\epsilon
  < \infty.
\]
Intuitively, this integral is finite when the class is not too
large: the number of brackets needed to cover $\mathcal{F}$ to
precision $\epsilon$ grows at most polynomially in $1/\epsilon$.

\begin{example}{Donsker and non-Donsker classes}
\label{w11:ex:donsker}
\leavevmode
\begin{itemize}
  \item \textbf{Parametric classes} (e.g., logistic regression
    models indexed by a finite-dimensional parameter) are Donsker
    under mild moment conditions, since the bracketing entropy
    grows only logarithmically.
  \item \textbf{Hölder-smooth function classes} on $[0,1]^d$
    with smoothness index $s > d/2$ are Donsker; the entropy
    integral converges because the bracketing numbers grow
    polynomially in $1/\epsilon$.
  \item \textbf{Indicator classes} $\mathcal{F} = \{x \mapsto
    \mathbf{1}(x \leq t) : t \in \mathbb{R}\}$ are Donsker
    by the classical Donsker theorem for empirical distribution
    functions.
  \item \textbf{Random forests and neural networks} typically fall
    outside classical Donsker regimes in the forms used for modern
    flexible prediction, because the function classes they implicitly
    define can approximate arbitrarily complex functions.
  \item \textbf{High-dimensional lasso} with a growing number of
    selected variables also typically falls outside Donsker regimes,
    because the effective model dimension grows with $n$.
\end{itemize}
\end{example}

\subsection{Connecting the two terms}
\label{w11:subsec:donsker-role}

Returning to the remainder \eqref{w11:eq:Rn-decomp}, the two terms
correspond exactly to the two challenges identified in
Section~\ref{w11:subsec:infinite-dim}:
\begin{itemize}
  \item \textbf{Population bias.}  Controlled by Neyman
    orthogonality alone; $o_p(n^{-1/2})$ under the product-rate
    condition regardless of how $\hat\eta$ is estimated.
  \item \textbf{Empirical process term.}  Requires a separate
    argument.  If $\{f_\eta : \eta \in \mathcal{H}\}$ is Donsker
    and $\|\hat\eta - \eta_0\| \to 0$ in probability, then by
    stochastic equicontinuity,
    \begin{align}
      \mathbb{G}_n\{f_{\hat\eta} - f_{\eta_0}\} = o_p(1),
      \label{w11:eq:donsker-control}
    \end{align}
    making the empirical process term $o_p(n^{-1/2})$.
    Combined with the controlled population bias, the full
    remainder $R_n = o_p(n^{-1/2})$, and root-$n$ asymptotic
    linearity follows from \eqref{w11:eq:plugin-expand}.
\end{itemize}
When the Donsker condition fails, \eqref{w11:eq:donsker-control}
need not hold: the empirical process term can remain non-negligible
even as $\hat\eta \to \eta_0$.  Root-$n$ asymptotic linearity of
the plug-in estimator then breaks down even though the score is
orthogonal and the nuisance estimators are consistent.  This is why
orthogonality and cross-fitting are complements rather than
substitutes: orthogonality handles Problem~1 and cross-fitting
handles Problem~2, and both are needed for valid root-$n$ inference
with flexible nuisance learners.

\begin{remark}[Donsker conditions and machine learning]
\label{w11:rem:donsker}
The Donsker condition fails for random forests, neural networks,
and high-dimensional lasso precisely because their function classes
are too rich to have finite bracketing entropy.  This is a genuine
obstruction when flexible learners are used without data splitting.
Cross-fitting sidesteps it entirely: because $\hat\eta^{(-k)}$ is
independent of the observations in $\mathcal{I}_k$, the empirical
process term \eqref{w11:eq:ep-term} is replaced by a sum of
conditionally independent terms, each $o_p(n^{-1/2})$ by a
conditional law of large numbers, with no Donsker assumption
required.
\end{remark}

The practical solution is therefore to separate nuisance estimation
from score evaluation, which leads to sample splitting and, in its
more efficient form, cross-fitting.


% ============================================================
\section{Sample Splitting}
\label{w11:sec:splitting}
% ============================================================

Section~\ref{w11:sec:reuse} showed that the plug-in remainder $R_n$
in \eqref{w11:eq:Rn-decomp} has two parts: the population bias,
handled by Neyman orthogonality, and the empirical process term
\eqref{w11:eq:ep-term}, which is not controlled by orthogonality
alone when the Donsker condition fails.  Sample splitting resolves
the second problem by construction.

Partition the indices $\{1,\ldots,n\}$ into a training sample
$\mathcal{I}_{\mathrm{train}}$ and an evaluation sample
$\mathcal{I}_{\mathrm{eval}}$.  Fit the nuisance estimators
\[
  \hat\eta^{(\mathrm{train})}
  =
  \bigl(
    \hat\mu_0^{(\mathrm{train})},\;
    \hat\mu_1^{(\mathrm{train})},\;
    \hat\pi^{(\mathrm{train})}
  \bigr)
\]
on $\mathcal{I}_{\mathrm{train}}$, then define $\hat\tau$ by
solving the estimating equation on the held-out sample:
\[
  \frac{1}{|\mathcal{I}_{\mathrm{eval}}|}
  \sum_{i \in \mathcal{I}_{\mathrm{eval}}}
  \varphi_{\mathrm{eff}}\!\bigl(O_i;\,\tau,\,\hat\eta^{(\mathrm{train})}\bigr)
  = 0.
\]

To see why this eliminates the empirical process problem, write the
analog of the remainder decomposition \eqref{w11:eq:Rn-decomp} on
the evaluation sample.  Let $n_{\mathrm{e}} = |\mathcal{I}_{\mathrm{eval}}|$
and $\mathbb{P}_{n_{\mathrm{e}}}$ denote the empirical measure over
$\mathcal{I}_{\mathrm{eval}}$.  The remainder on the evaluation fold is
\begin{align}
  R_n^{(\mathrm{e})}
  &=
  \underbrace{%
    P\{
      \varphi_{\mathrm{eff}}(\,\cdot\,;\tau,\hat\eta^{(\mathrm{train})})
      - \varphi_{\mathrm{eff}}(\,\cdot\,;\tau,\eta_0)
    \}
  }_{\text{population bias: } o_p(n^{-1/2}) \text{ by orthogonality}}
  \nonumber\\[4pt]
  &\quad+
  \underbrace{%
    (\mathbb{P}_{n_{\mathrm{e}}} - P)\{
      \varphi_{\mathrm{eff}}(\,\cdot\,;\tau,\hat\eta^{(\mathrm{train})})
      - \varphi_{\mathrm{eff}}(\,\cdot\,;\tau,\eta_0)
    \}
  }_{\text{empirical process term}}.
  \label{w11:eq:split-decomp}
\end{align}
The population bias is $o_p(n^{-1/2})$ by Neyman orthogonality,
exactly as before.  The empirical process term is now different in
character: because $\hat\eta^{(\mathrm{train})}$ is estimated on
$\mathcal{I}_{\mathrm{train}}$, it is \emph{independent} of the
observations $\{O_i : i \in \mathcal{I}_{\mathrm{eval}}\}$.
Conditioning on $\hat\eta^{(\mathrm{train})}$, the summands in the
empirical process term are independent and mean zero, so by the law
of large numbers applied conditionally,
\[
  (\mathbb{P}_{n_{\mathrm{e}}} - P)\{
    \varphi_{\mathrm{eff}}(\,\cdot\,;\tau,\hat\eta^{(\mathrm{train})})
    - \varphi_{\mathrm{eff}}(\,\cdot\,;\tau,\eta_0)
  \}
  = o_p(n^{-1/2}).
\]
No Donsker condition on the function class is needed: independence
does the work that Donsker equicontinuity would otherwise have to do.
The full remainder $R_n^{(\mathrm{e})} = o_p(n^{-1/2})$, and
root-$n$ asymptotic linearity of the split estimator follows.

The main disadvantage of sample splitting is inefficiency: only
$|\mathcal{I}_{\mathrm{eval}}|$ observations contribute to the
estimation of $\tau$, and the estimate depends on the particular
random partition chosen.  In other words, sample splitting buys
cleaner asymptotic theory by sacrificing training sample size for
each nuisance learner, which can slow the convergence of $\hat\eta$
and inflate the variance of the final estimator.

\begin{remark}
\label{w11:rem:split-noise}
A single split can be noisy.  In practice, sensitivity to the split
is often assessed by repeating the procedure with several random
partitions and aggregating results.  Cross-fitting provides a more
systematic solution by rotating the training and evaluation roles
across all observations.
\end{remark}


% ============================================================
\section{Cross-Fitting}
\label{w11:sec:crossfit}
% ============================================================

Cross-fitting extends sample splitting to recover full-sample
efficiency while preserving the independence argument that
eliminates the empirical process term.  The key observation from
Section~\ref{w11:sec:splitting} was that conditioning on
$\hat\eta^{(\mathrm{train})}$ makes the evaluation-fold summands
independent and mean zero, so the empirical process term vanishes
by a conditional law of large numbers.  Cross-fitting applies this
argument simultaneously to every observation by rotating the
training and evaluation roles across $K$ folds.

Partition the indices $\{1,\ldots,n\}$ into $K$ approximately
equal folds $\mathcal{I}_1,\ldots,\mathcal{I}_K$.  For each
fold $k$:
\begin{enumerate}
  \item Estimate the nuisance functions using all observations
    \emph{not} in fold $k$, obtaining
    \[
      \hat\eta^{(-k)}
      =
      \bigl(
        \hat\mu_0^{(-k)},\;
        \hat\mu_1^{(-k)},\;
        \hat\pi^{(-k)}
      \bigr).
    \]
  \item Evaluate the orthogonal score on the held-out fold
    $\mathcal{I}_k$ using $\hat\eta^{(-k)}$.
  \item Aggregate the score contributions across all $K$ folds.
\end{enumerate}
The cross-fitted estimator $\hat\tau_{\mathrm{DML}}$ solves
\begin{align}
  \frac{1}{n}
  \sum_{k=1}^K
  \sum_{i \in \mathcal{I}_k}
  \varphi_{\mathrm{eff}}\!\bigl(O_i;\,\tau,\,\hat\eta^{(-k)}\bigr)
  = 0.
  \label{w11:eq:crossfit-moment}
\end{align}

To see why the empirical process problem is resolved, write the
remainder \eqref{w11:eq:Rn-decomp} as a sum over folds:
\[
  R_n = \frac{1}{K}\sum_{k=1}^K R_n^{(k)},
\]
where $R_n^{(k)}$ is the remainder contribution from fold $k$,
decomposed exactly as in \eqref{w11:eq:split-decomp} with
$\mathcal{I}_{\mathrm{eval}} = \mathcal{I}_k$ and
$\hat\eta^{(\mathrm{train})} = \hat\eta^{(-k)}$.  For each $k$,
the nuisance estimate $\hat\eta^{(-k)}$ is trained on
$\{1,\ldots,n\}\setminus\mathcal{I}_k$, which is independent of
the observations in $\mathcal{I}_k$.  Conditioning on
$\hat\eta^{(-k)}$, the fold-$k$ summands are independent and mean
zero, and the conditional law of large numbers gives
$R_n^{(k)} = o_p(n^{-1/2})$ for each $k$ without any Donsker
assumption.  Averaging over folds, $R_n = o_p(n^{-1/2})$.

Cross-fitting therefore achieves two goals simultaneously: it
preserves the held-out independence structure that removes the need
for Donsker conditions, and it ensures every observation serves
once as an evaluation point, so the full sample contributes to the
estimation of $\tau$.  For this reason, cross-fitting has become
the standard device for combining orthogonal scores with flexible
nuisance estimation \citep{chernozhukov2018double,vanderLaan2011targeted}.
Cross-fitting removes the same-sample empirical-process obstruction,
but it does not by itself guarantee good finite-sample performance
if the nuisance estimators are inaccurate or overlap is poor.

\begin{tcolorbox}[defstyle,
  title={Algorithm: Cross-Fitted Orthogonal-Score Estimator for the ATE}]
\label{w11:alg:crossfit}
\begin{enumerate}
  \item Partition the sample into $K$ folds
    $\mathcal{I}_1,\ldots,\mathcal{I}_K$.
  \item For each $k = 1,\ldots,K$, fit nuisance estimators
    $\hat\eta^{(-k)} = (\hat\mu_0^{(-k)}, \hat\mu_1^{(-k)},
    \hat\pi^{(-k)})$ on the complement
    $\{1,\ldots,n\} \setminus \mathcal{I}_k$.
  \item For each observation $i \in \mathcal{I}_k$, compute the
    out-of-fold score contribution
    \[
      \hat\varphi_i
      =
      \varphi_{\mathrm{eff}}\!\bigl(O_i;\,\tau,\,\hat\eta^{(-k)}\bigr).
    \]
  \item Solve the cross-fitted moment equation
    \eqref{w11:eq:crossfit-moment}: because the score is linear
    in $\tau$, the solution is
    \begin{align}
      \hat\tau_{\mathrm{DML}} = \hat\mu_{1,\mathrm{DML}} - \hat\mu_{0,\mathrm{DML}},
      \label{w11:eq:dml-ate}
    \end{align}
    where
    \begin{align*}
      \hat\mu_{1,\mathrm{DML}}
      &=
      \frac{1}{n}\sum_{k=1}^{K}\sum_{i\in\mathcal{I}_k}
      \Bigl[
        \hat\mu_1^{(-k)}(X_i)
        +
        \frac{T_i}{\hat\pi^{(-k)}(X_i)}
        \bigl\{Y_i - \hat\mu_1^{(-k)}(X_i)\bigr\}
      \Bigr], \\[6pt]
      \hat\mu_{0,\mathrm{DML}}
      &=
      \frac{1}{n}\sum_{k=1}^{K}\sum_{i\in\mathcal{I}_k}
      \Bigl[
        \hat\mu_0^{(-k)}(X_i)
        +
        \frac{1-T_i}{1-\hat\pi^{(-k)}(X_i)}
        \bigl\{Y_i - \hat\mu_0^{(-k)}(X_i)\bigr\}
      \Bigr].
    \end{align*}
\end{enumerate}
\end{tcolorbox}


% ============================================================
\section{Double Machine Learning for the Average Treatment Effect}
\label{w11:sec:dml}
% ============================================================

The estimator in \eqref{w11:eq:dml-ate} is simply the AIPW estimator
from Chapter~10, \eqref{w10:eq:aipw}, with out-of-fold nuisance
estimates substituted in place of estimates fit on the full sample.
The term \emph{double machine learning} (DML), introduced by
\citet{chernozhukov2018double}, refers to the combination of three
ingredients:
\begin{itemize}
  \item \emph{Neyman orthogonality}, which eliminates the
    population bias term in the remainder decomposition
    \eqref{w11:eq:Rn-decomp}, so nuisance estimation error
    enters only at second order;
  \item \emph{cross-fitting}, which eliminates the empirical
    process term in \eqref{w11:eq:Rn-decomp} via fold-by-fold
    conditional independence, removing the need for Donsker
    conditions; and
  \item \emph{flexible machine-learning estimation} of the
    nuisance functions $\mu_0$, $\mu_1$, and $\pi$, which is
    made valid by the first two ingredients.
\end{itemize}
The word ``double'' should not be read as merely meaning ``two
models''; it refers to the combination of nuisance learning with
orthogonalization and debiasing.  The crucial feature is the
orthogonal moment structure combined with cross-fitting, which
together make flexible nuisance estimation compatible with root-$n$
inference.  In this sense, DML is not a new causal estimand and not
a new identification strategy; it is a modern estimation protocol
built around an orthogonal score.

\begin{remark}
\label{w11:rem:cf-aipw}
The estimator in \eqref{w11:eq:dml-ate} is often called the
\emph{cross-fitted AIPW estimator} in the biostatistics literature.
The two names refer to the same object and may be used
interchangeably when the context is clear.
\end{remark}

\begin{warn}{``Double'' does not mean ``two models are enough''}
\label{w11:warn:double}
It is tempting to think that estimating any two of $\mu_0$,
$\mu_1$, and $\pi$ flexibly is sufficient.  What matters is that
all nuisance components entering the orthogonal score are estimated
out-of-fold, and that the product-rate condition
(Section~\ref{w11:sec:rates}) is satisfied.  Using flexible methods
for only one nuisance function while misspecifying another can still
produce invalid inference.
\end{warn}


% ============================================================
\section{Rate Conditions and Asymptotic Normality}
\label{w11:sec:rates}
% ============================================================

The main theoretical advantage of orthogonality and cross-fitting is
that valid root-$n$ inference can hold even when nuisance estimators
converge more slowly than $n^{-1/2}$.  The following result makes
this precise.

\begin{thm}{Asymptotic Linearity Under Cross-Fitting}
\label{w11:thm:asymp}
Suppose:
\begin{enumerate}[label=(\roman*)]
  \item the score $\varphi_{\mathrm{eff}}(O;\,\tau,\eta)$ is
    orthogonal at the truth $(\tau_0, \eta_0)$;
  \item the nuisance estimators are consistent in mean squared error;
  \item \emph{strong overlap} holds: there exists a constant $c > 0$
    such that
    \[
      c \;\leq\; \pi(X) \;\leq\; 1 - c
      \quad\text{and}\quad
      c \;\leq\; \hat\pi^{(-k)}(X) \;\leq\; 1 - c
    \]
    almost surely for each fold $k$;
  \item the nuisance estimators satisfy the
    \emph{product-rate condition}
    \begin{align}
      \|\hat\pi^{(-k)} - \pi\| \cdot
      \|\hat\mu_t^{(-k)} - \mu_t\|
      = o_p(n^{-1/2}),
      \qquad t = 0, 1,
      \label{w11:eq:product-rate}
    \end{align}
    where $\|\cdot\|$ denotes the $L_2(P)$ norm.
\end{enumerate}
Then the cross-fitted estimator satisfies
\[
  \sqrt{n}(\hat\tau_{\mathrm{DML}} - \tau)
  =
  \frac{1}{\sqrt{n}}\sum_{i=1}^n
  \varphi_{\mathrm{eff}}(O_i;\,\tau,\eta)
  + o_p(1),
\]
and consequently
\[
  \sqrt{n}(\hat\tau_{\mathrm{DML}} - \tau)
  \overset{d}{\longrightarrow}
  N\!\Bigl(0,\;\E\bigl[\varphi_{\mathrm{eff}}(O;\,\tau,\eta)^2\bigr]\Bigr).
\]
\end{thm}

The asymptotic variance equals the semiparametric efficiency bound
(Theorem~\ref{w10:thm:bound}), so $\hat\tau_{\mathrm{DML}}$ is
asymptotically efficient whenever the product-rate condition holds.

\begin{remark}[Weak Overlap vs.\ Strong Overlap]
\label{w11:rem:overlap}
Condition~(iii) in the theorem requires \emph{strong overlap}: both
the true propensity score $\pi(X)$ and the estimated propensity score
$\hat\pi^{(-k)}(X)$ must be bounded uniformly away from $0$ and $1$
by a positive constant $c$.  It is important to distinguish this from
the weaker positivity condition $0 < \pi(X) < 1$ a.s., which serves
a different role.

\emph{Weak overlap} ($0 < \pi(X) < 1$ a.s.) is the condition needed
for \emph{identification}: it ensures that every covariate stratum
contains units from both treatment arms, so that the counterfactual
means $\E[Y(t) \mid X]$ are identified from observed data.  Without
weak overlap, certain strata have no counterfactual information and the
ATE is not nonparametrically identified.

\emph{Strong overlap} ($c \leq \pi(X) \leq 1-c$ and
$c \leq \hat\pi^{(-k)}(X) \leq 1-c$ for some $c > 0$) is what
the proof above actually uses: it bounds the inverse-probability
weights $1/\hat\pi$ and $1/(1-\hat\pi)$ away from infinity, ensuring
that the Cauchy-Schwarz step in Step~1 and the variance bound in
Step~2 go through.  Without strong overlap, the influence function
values can have infinite variance, and $\sqrt{n}$-consistent,
asymptotically normal inference breaks down even if the ATE is
identified \citep{khan2010irregular}.

The practical consequence is that weak overlap is necessary but not
sufficient for regular inference.  When the propensity score
approaches 0 or 1 in parts of the covariate space, the inverse
weights blow up, influence function values become extreme, and the
variance estimator from Theorem~\ref{w11:thm:var} can be severely
unstable.  Trimming extreme weights or restricting the target
population to a region of stronger overlap are common practical
responses.  The requirement on $\hat\pi^{(-k)}$ is equally
important: even if the true $\pi$ is well-behaved, an estimated
propensity score that strays near $0$ or $1$ in a given fold will
cause the same instability through the finite-sample weights.
\end{remark}

\begin{proof}[Proof sketch]
Write $n_k = |\mathcal{I}_k|$ and $\mathbb{P}_{n_k}$ for the
empirical measure over fold $k$.  The cross-fitted moment equation
\eqref{w11:eq:crossfit-moment} gives
\[
  0 = \frac{1}{K}\sum_{k=1}^K \mathbb{P}_{n_k}
      \bigl\{\varphi_{\mathrm{eff}}(O;\,\tau,\hat\eta^{(-k)})\bigr\}.
\]
Adding and subtracting $\varphi_{\mathrm{eff}}(O;\,\tau,\eta_0)$
inside each fold and rearranging, using the fact that the score is
linear in $\tau$ with $\partial_\tau\varphi_{\mathrm{eff}} = -1$,
\[
  \hat\tau_{\mathrm{DML}} - \tau
  =
  \frac{1}{n}\sum_{i=1}^n \varphi_{\mathrm{eff}}(O_i;\,\tau,\eta_0)
  +
  \frac{1}{K}\sum_{k=1}^K R_n^{(k)},
\]
where each fold remainder $R_n^{(k)}$ decomposes as in
\eqref{w11:eq:split-decomp}:
\[
  R_n^{(k)}
  =
  \underbrace{%
    P\bigl\{\varphi_{\mathrm{eff}}(\,\cdot\,;\tau,\hat\eta^{(-k)})
           - \varphi_{\mathrm{eff}}(\,\cdot\,;\tau,\eta_0)\bigr\}
  }_{\text{population bias } B_n^{(k)}}
  +
  \underbrace{%
    (\mathbb{P}_{n_k} - P)\bigl\{\varphi_{\mathrm{eff}}(\,\cdot\,;\tau,\hat\eta^{(-k)})
           - \varphi_{\mathrm{eff}}(\,\cdot\,;\tau,\eta_0)\bigr\}
  }_{\text{empirical process term } E_n^{(k)}}.
\]
It suffices to show $\sqrt{n}\,R_n^{(k)} = o_p(1)$ for each $k$.

\medskip
\noindent\textit{Step 1: population bias.}\enspace
By iterated expectations conditioning on $X$, and using
$\E(T\mid X) = \pi(X)$ and $\E(Y\mid T{=}t, X) = \mu_t(X)$,
\begin{align*}
  B_n^{(k)}
  &= \E\!\left[
      \frac{\pi(X)}{\hat\pi^{(-k)}(X)}\bigl\{\mu_1(X)-\hat\mu_1^{(-k)}(X)\bigr\}
      -
      \frac{1-\pi(X)}{1-\hat\pi^{(-k)}(X)}\bigl\{\mu_0(X)-\hat\mu_0^{(-k)}(X)\bigr\}
    \right.\\
  &\qquad\qquad\left.
      +\,\hat\mu_1^{(-k)}(X)-\mu_1(X)
      -\,\hat\mu_0^{(-k)}(X)+\mu_0(X)
    \right]\\
  &= -\E\!\left[
      \bigl(\hat\mu_1^{(-k)}-\mu_1\bigr)(X)
      \cdot
      \frac{\hat\pi^{(-k)}(X)-\pi(X)}{\hat\pi^{(-k)}(X)}
      -
      \bigl(\hat\mu_0^{(-k)}-\mu_0\bigr)(X)
      \cdot
      \frac{\hat\pi^{(-k)}(X)-\pi(X)}{1-\hat\pi^{(-k)}(X)}
    \right].
\end{align*}
Under strong overlap (condition~(iii)), $\hat\pi^{(-k)}$ is bounded
away from $0$ and $1$ with probability approaching one, so the
denominators are bounded below.  Applying the Cauchy-Schwarz
inequality,
\[
  |B_n^{(k)}|
  \leq
  C\Bigl(
    \|\hat\mu_1^{(-k)} - \mu_1\|\cdot\|\hat\pi^{(-k)}-\pi\|
    +
    \|\hat\mu_0^{(-k)} - \mu_0\|\cdot\|\hat\pi^{(-k)}-\pi\|
  \Bigr)
  = o_p(n^{-1/2})
\]
by the product-rate condition~(iv).  Hence $\sqrt{n}\,B_n^{(k)} =
o_p(1)$.

\medskip
\noindent\textit{Step 2: empirical process term.}\enspace
Condition on $\hat\eta^{(-k)}$.  Since $\hat\eta^{(-k)}$ is
trained on $\{1,\ldots,n\}\setminus\mathcal{I}_k$, the summands
$\bigl\{\varphi_{\mathrm{eff}}(O_i;\tau,\hat\eta^{(-k)}) -
\varphi_{\mathrm{eff}}(O_i;\tau,\eta_0)\bigr\}$ for
$i\in\mathcal{I}_k$ are conditionally i.i.d.\ with conditional mean
zero (the mean is $B_n^{(k)}$, which has already been removed).
Let $f_k = \varphi_{\mathrm{eff}}(\cdot;\tau,\hat\eta^{(-k)}) -
\varphi_{\mathrm{eff}}(\cdot;\tau,\eta_0)$.  The conditional
variance of $E_n^{(k)}$ satisfies
\[
  \mathrm{Var}\bigl(E_n^{(k)}\mid\hat\eta^{(-k)}\bigr)
  =
  \frac{1}{n_k}\,\mathrm{Var}\bigl(f_k(O)\mid\hat\eta^{(-k)}\bigr)
  \leq
  \frac{1}{n_k}\,\|f_k\|_{L_2}^2.
\]
Under strong overlap and consistency (conditions~(ii)--(iii)), the
$L_2$ norm satisfies $\|f_k\|_{L_2}^2 = o_p(1)$, so
\[
  \mathrm{Var}\bigl(\sqrt{n}\,E_n^{(k)}\mid\hat\eta^{(-k)}\bigr)
  =
  \frac{n}{n_k}\,\|f_k\|_{L_2}^2
  = O(1)\cdot o_p(1)
  = o_p(1),
\]
using $n/n_k = O(1)$ since folds are of equal size.  By conditional
Chebyshev, $\sqrt{n}\,E_n^{(k)} = o_p(1)$.

\medskip
\noindent\textit{Conclusion.}\enspace
Combining Steps~1 and~2, $\sqrt{n}\,R_n^{(k)} = o_p(1)$ for each
$k$, and hence $\sqrt{n}\,\frac{1}{K}\sum_k R_n^{(k)} = o_p(1)$.
The first term $n^{-1/2}\sum_i\varphi_{\mathrm{eff}}(O_i;\tau,\eta_0)$
converges in distribution to $N(0, \E[\varphi_{\mathrm{eff}}^2])$
by the ordinary CLT.  Asymptotic linearity and asymptotic normality
follow by Slutsky's theorem.
\end{proof}

\begin{remark}[Interpreting the product-rate condition]
\label{w11:rem:product-rate}
Condition~\eqref{w11:eq:product-rate} requires the \emph{product}
of the two nuisance estimation errors to be $o_p(n^{-1/2})$.  This
is substantially weaker than requiring each error individually to
be $o_p(n^{-1/2})$.  A sufficient symmetric case is that each nuisance estimator
converges at least at rate $o_p(n^{-1/4})$ in mean square —
achievable by many nonparametric and machine-learning methods
under regularity conditions — in which case the product is
$o_p(n^{-1/2})$, satisfying the condition.  This is what makes
flexible nuisance estimation feasible in semiparametric causal
inference.
\end{remark}

\begin{remark}[Connection to Chapter~10]
\label{w11:rem:ch10-link}
Theorem~\ref{w11:thm:asymp} is the cross-fitting version of the
asymptotic normality result stated in Chapter~10,
Section~\ref{w10:sec:asymp}.  The product-rate condition
\eqref{w11:eq:product-rate} is identical to \eqref{w10:eq:product-rate};
the difference is that here the nuisance estimates are out-of-fold,
which removes the need for Donsker conditions.  Cross-fitting thus
provides the same asymptotic conclusion while accommodating a
substantially broader class of nuisance learners.
\end{remark}



% ============================================================
\section{Lab: Simulation Study of the DML Estimator}
\label{w11:sec:lab}
% ============================================================

This lab uses simulation to verify the theoretical properties of the
DML estimator established in Theorem~\ref{w11:thm:asymp}.  The
central message of this chapter is that flexible nuisance estimation
alone is not sufficient for valid inference: even a consistent
nonparametric estimator can produce a biased AIPW estimator if the
same sample is reused for nuisance estimation and score evaluation.
Cross-fitting resolves this by eliminating the empirical process
term in the remainder decomposition~\eqref{w11:eq:Rn-decomp}.

To isolate this data-reuse problem from model misspecification, all
three estimators below use \emph{random forests} for nuisance
estimation.  The only difference between the naive and DML estimators
is whether the nuisance estimates are produced on the full sample
(same data reused) or on held-out folds (cross-fitted).

\subsection{Data-generating process}
\label{w11:subsec:lab-dgp}

Let $n = 2000$ and generate covariates $X_1, X_2, X_3
\overset{\mathrm{iid}}{\sim} \mathrm{Uniform}(-1, 1)$, collected as
$X = (X_1, X_2, X_3)^\top$.  Define the true nuisance functions by
\begin{align*}
  \pi(X)
  &= \frac{\exp\!\bigl(\sin(\pi X_1) + X_2^2 - 0.5\bigr)}
          {1 + \exp\!\bigl(\sin(\pi X_1) + X_2^2 - 0.5\bigr)}, \\[4pt]
  \mu_1(X) &= 2\sin(\pi X_1) + X_2^2 + X_3, \\[4pt]
  \mu_0(X) &= \sin(\pi X_1) + X_2^2 - X_3.
\end{align*}
The true ATE is
\[
  \tau = \E\{\mu_1(X) - \mu_0(X)\} = \E\{\sin(\pi X_1) + 2X_3\} = 0,
\]
since $X_1$ and $X_3$ are symmetric about zero.  Treatment and
outcome are drawn as
\[
  T_i \mid X_i \sim \mathrm{Bernoulli}\{\pi(X_i)\},
  \qquad
  Y_i = \mu_{T_i}(X_i) + \varepsilon_i,
  \qquad
  \varepsilon_i \sim N(0, 1).
\]
The nonlinearity of $\pi$ and $\mu_t$ means that random forests can
fit them consistently, while the Donsker condition fails for the
function class they trace out — exactly the setting of
Section~\ref{w11:sec:reuse}.

\subsection{The three estimators}
\label{w11:subsec:lab-estimators}

\begin{enumerate}
  \item \textbf{Oracle AIPW.}  Plug in the true nuisance functions
    $\pi(X_i)$, $\mu_1(X_i)$, $\mu_0(X_i)$ directly.  This
    estimator is infeasible but provides the semiparametric
    efficiency benchmark.
  \item \textbf{Naive AIPW.}  Estimate $\pi$, $\mu_1$, $\mu_0$
    using random forests on the \emph{full sample}, then evaluate
    the AIPW score on the same sample.  The nuisance estimators are
    consistent, but data reuse means the empirical process term in
    \eqref{w11:eq:Rn-decomp} is not controlled — the Donsker
    condition fails for random forests
    (Remark~\ref{w11:rem:donsker}).
  \item \textbf{DML (cross-fitted AIPW).}  Estimate $\pi$, $\mu_1$,
    $\mu_0$ using the same random forest learners, but via $K = 5$
    fold cross-fitting.  The fold-by-fold independence argument of
    Section~\ref{w11:sec:crossfit} eliminates the empirical process
    term without any Donsker assumption.
\end{enumerate}

Estimators (2) and (3) use identical learners; the only difference
is data reuse versus cross-fitting.

\subsection{Simulation results and interpretation}
\label{w11:subsec:lab-results}

Running the simulation with $n = 2000$ and $n_{\mathrm{sim}} = 500$
replications produces the following results:

\medskip
\begin{center}
\renewcommand{\arraystretch}{1.3}
\begin{tabular}{lccc}
\hline
\textbf{Estimator} & \textbf{Bias} & \textbf{Variance} & \textbf{RMSE} \\
\hline
Oracle AIPW   & $-0.003$ & $0.003$ & $0.056$ \\
Naive RF AIPW & $\phantom{-}0.120$ & $0.006$ & $0.142$ \\
DML           & $\phantom{-}0.077$ & $0.015$ & $0.145$ \\
\hline
\end{tabular}
\end{center}
\medskip

The simulation illustrates the asymptotic logic of cross-fitting:
it reduces the bias created by same-sample nuisance fitting, but
finite-sample performance still depends on the quality of the
nuisance learners, the degree of overlap, and the sample size.

\begin{itemize}
  \item \textbf{Oracle AIPW} is nearly unbiased with small variance,
    serving as the semiparametric efficiency benchmark.
  \item \textbf{Naive RF AIPW} exhibits substantial bias (0.120)
    despite using consistent random forest estimators.  The bias
    arises from the uncontrolled empirical process term
    \eqref{w11:eq:ep-term}: reusing the same sample for nuisance
    estimation and score evaluation allows overfitting in the
    random forests to propagate into the AIPW score, exactly as
    analyzed in Section~\ref{w11:sec:reuse}.
  \item \textbf{DML} reduces the bias to 0.077 and, as predicted by
    theory, increases the variance (from 0.006 to 0.015).
    Cross-fitting eliminates the empirical process term via the
    fold-by-fold independence argument of
    Section~\ref{w11:sec:crossfit}, which accounts for the
    bias reduction.  The remaining bias reflects a finite-sample
    limitation: with $K = 5$ folds each training set contains only
    1600 observations, which may be insufficient for random forests
    to satisfy the product-rate condition
    \eqref{w11:eq:product-rate} at this sample size.  The condition
    is asymptotic, and convergence of DML bias to zero is expected
    as $n \to \infty$.
\end{itemize}

\begin{remark}
  The key comparison is between Naive RF AIPW and DML, not between
  either and the oracle.  Both use identical random forest learners;
  the only difference is cross-fitting.  The bias reduction from
  0.120 to 0.077 is therefore attributable solely to eliminating
  the data-reuse problem, while the residual bias of DML reflects
  the finite-sample gap between the actual nuisance convergence rate
  and the asymptotic product-rate condition.  Increasing $n$
  would close this gap.
\end{remark}

% ============================================================
\section{Variance Estimation and Confidence Intervals}
\label{w11:sec:variance}
% ============================================================

Inference for $\hat\tau_{\mathrm{DML}}$ proceeds exactly as in
Chapter~10, Section~\ref{w10:subsec:variance}, except that the
nuisance estimates are now out-of-fold.  Cross-fitting changes how
the nuisance functions are estimated, not the asymptotic variance
formula itself; the variance is still the variance of the efficient
score, estimated empirically using the out-of-fold score
contributions.  Let $k(i)$ denote the fold
containing observation $i$, and define the estimated influence value
\[
  \hat\varphi_i
  =
  \varphi_{\mathrm{eff}}\!\bigl(O_i;\,\hat\tau_{\mathrm{DML}},\,
  \hat\eta^{(-k(i))}\bigr).
\]
A natural variance estimator is
\begin{align}
  \hat{V}
  =
  \frac{1}{n(n-1)}\sum_{i=1}^n
  (\hat\varphi_i - \bar\varphi)^2,
  \qquad
  \bar\varphi = \frac{1}{n}\sum_{i=1}^n \hat\varphi_i.
  \label{w11:eq:var-est}
\end{align}
The Wald confidence interval at level $1 - \alpha$ is
\[
  \hat\tau_{\mathrm{DML}}
  \pm
  z_{1-\alpha/2}\sqrt{\hat{V}}.
\]

\begin{thm}{Consistency of the Variance Estimator}
\label{w11:thm:var}
Under the conditions of Theorem~\ref{w11:thm:asymp}, the variance
estimator $\hat{V}$ in \eqref{w11:eq:var-est} is consistent for
$\E[\varphi_{\mathrm{eff}}(O;\,\tau,\eta)^2]$, and the Wald
confidence interval is asymptotically valid.
\end{thm}

From a practical standpoint, the variance formula \eqref{w11:eq:var-est}
is simply the sample variance of the out-of-fold estimated influence
values.  The estimator is identical in form to the one in
Chapter~10; the only difference is that $\hat\varphi_i$ now uses
the out-of-fold nuisance estimate $\hat\eta^{(-k(i))}$ rather
than a full-sample nuisance estimate.

\begin{remark}
\label{w11:rem:var-centering}
The centering by $\bar\varphi$ in \eqref{w11:eq:var-est} is
asymptotically negligible because $\bar\varphi \to 0$ in
probability.  In small samples, centering is nonetheless
recommended as it prevents finite-sample inflation of the
variance estimate.
\end{remark}


% ============================================================
\section{A Practical Workflow}
\label{w11:sec:workflow}
% ============================================================

The theory above leads to the following workflow for applied
cross-fitted orthogonal-score estimation.  The steps are grouped
into three phases: identification and setup, estimation, and
diagnostics and interpretation.

\begin{tcolorbox}[defstyle, title={Workflow: Cross-Fitted Causal Estimation}]
\textbf{Phase 1: Identification and setup}
\begin{enumerate}
  \item \textbf{Specify the estimand.}  State the target causal
    parameter and verify the identification assumptions that justify
    it — conditional exchangeability, positivity, and consistency
    for the ATE (Chapter~5).
  \item \textbf{Choose the orthogonal score.}  For the ATE under
    conditional exchangeability, this is the efficient influence
    function \eqref{w11:eq:orth-score-ate}.  The score fixes what
    nuisance components must be estimated.
  \item \textbf{Select nuisance learners.}  Choose flexible learners
    for $\mu_0(x)$, $\mu_1(x)$, and $\pi(x)$.  Verify informally
    that the product-rate condition~\eqref{w11:eq:product-rate} is
    plausible: each learner should converge at rate at least
    $o_p(n^{-1/4})$ in mean square.
  \item \textbf{Monitor overlap.}  Before fitting, examine the
    empirical distribution of covariates by treatment group.
    Near-violation of positivity will inflate influence values and
    destabilize inference regardless of the estimator used.
\end{enumerate}

\medskip
\textbf{Phase 2: Estimation}
\begin{enumerate}[resume]
  \item \textbf{Partition and cross-fit.}  Partition the sample into
    $K$ folds ($K = 5$ is a common default).  For each fold $k$,
    fit the nuisance learners on the complement
    $\{1,\ldots,n\}\setminus\mathcal{I}_k$ to obtain out-of-fold
    estimates $\hat\eta^{(-k)}$.
  \item \textbf{Compute out-of-fold orthogonal scores.}  For each
    $i \in \mathcal{I}_k$, compute
    $\hat\varphi_i =
    \varphi_{\mathrm{eff}}(O_i;\,\hat\tau_{\mathrm{DML}},\,
    \hat\eta^{(-k)})$.
  \item \textbf{Solve the moment equation.}  Average the scores
    across all observations to obtain $\hat\tau_{\mathrm{DML}}$
    via the explicit formula~\eqref{w11:eq:dml-ate}.
  \item \textbf{Estimate the variance.}  Compute $\hat{V}$ as the
    empirical variance of the out-of-fold scores
    \eqref{w11:eq:var-est} and form the Wald confidence interval.
\end{enumerate}

\medskip
\textbf{Phase 3: Diagnostics and interpretation}
\begin{enumerate}[resume]
  \item \textbf{Inspect estimated propensity scores.}  Examine the
    distribution of $\hat\pi^{(-k)}(X_i)$ across folds; extreme
    values near $0$ or $1$ inflate influence-value variance and can
    destabilize coverage.
  \item \textbf{Assess sensitivity.}  Repeat the analysis with
    alternative learners for the nuisance functions.  Substantial
    sensitivity to learner choice signals that the product-rate
    condition may not be satisfied at this sample size.
  \item \textbf{Interpret relative to identification assumptions.}
    The credibility of the causal estimate rests on the
    identification assumptions, not on the predictive accuracy of
    the nuisance learners.  These assumptions are untestable from
    the data and must be argued on substantive grounds.
\end{enumerate}
\end{tcolorbox}

Step~11 deserves emphasis.  A common error is to assess the
credibility of a causal estimate by citing good out-of-sample
predictive accuracy for the nuisance models.  Predictive accuracy
is necessary but not sufficient.  The identifying assumptions ---
conditional exchangeability, positivity, consistency --- are
untestable from the data and must be argued on substantive grounds.

\begin{quote}
  Modern causal estimation is not merely machine learning plus a
  causal estimand.  It is machine learning embedded inside a
  carefully constructed semiparametric procedure, whose validity
  rests on identification assumptions that no algorithm can verify.
\end{quote}


% ============================================================
\section{What Machine Learning Does Not Solve}
\label{w11:sec:limits}
% ============================================================

Flexible nuisance estimation is a genuine improvement over
parametric methods, but it does not resolve the fundamental
difficulties of causal inference.  In particular, machine learning
does not solve:
\begin{itemize}
  \item \textbf{Unmeasured confounding.}  If important confounders
    are absent from $X$, conditional exchangeability
    $Y(t) \indep T \mid X$ fails.  No amount of flexibility in
    estimating $\pi(x)$ or $\mu_t(x)$ can correct this.
  \item \textbf{Violations of consistency.}  If the treatment version
    is ill-defined or the stable unit treatment value assumption
    (SUTVA) is violated, the potential outcomes framework itself
    is compromised.
  \item \textbf{Failure of positivity.}  When $\pi(x) \approx 0$
    or $\pi(x) \approx 1$ for part of the covariate space,
    influence values become unstable and coverage of confidence
    intervals can degrade severely.
  \item \textbf{Ambiguity about the estimand.}  Flexible estimation
    cannot resolve disagreement about what causal quantity is
    scientifically meaningful.
  \item \textbf{Invalid instrumental variable assumptions.}
    The exclusion restriction and exogeneity conditions
    (Chapter~7) are not testable from the data; machine learning
    cannot substitute for subject-matter justification.
  \item \textbf{Fragile mediation assumptions.}
    Sequential ignorability and no interference between
    mediators (Chapter~8) must be argued on substantive grounds.
\end{itemize}

Machine learning improves the estimation of nuisance functions
\emph{within} a given identification strategy.  It does not create
identification where none exists.  Machine learning can improve
estimation after the causal problem has been formulated and
identified, but it cannot repair a failure of identification
assumptions.

\begin{remark}
\label{w11:rem:ml-limits}
Students sometimes encounter the claim that machine learning can
``solve'' causal inference problems.  The identification step
remains entirely the responsibility of the researcher; no
algorithm can substitute for subject-matter justification of the
core assumptions.
\end{remark}



% ============================================================
\section{Chapter Summary}
\label{w11:sec:summary}
% ============================================================

\begin{enumerate}
  \item Flexible methods are attractive for nuisance estimation when
    the true nuisance functions are complex, but naive plug-in
    estimators can fail for two reasons: orthogonality may be
    absent, and even when orthogonality holds, same-sample nuisance
    fitting can create an additional empirical-process bias
    (Section~\ref{w11:sec:plugin}).
  \item A score function is Neyman orthogonal if its Gateaux
    derivative with respect to the nuisance vanishes at the truth,
    so nuisance estimation error enters only at second order
    (Definition~\ref{w11:defn:orthogonal}).
  \item The efficient influence function for the ATE
    \eqref{w11:eq:orth-score-ate} is orthogonal, doubly robust,
    and efficiency-achieving.  It serves as the orthogonal score
    throughout this chapter (Section~\ref{w11:sec:ate-score}).
  \item Even with an orthogonal score, reusing the same data for
    nuisance estimation and score evaluation can induce overfitting
    bias, especially when Donsker conditions fail
    (Section~\ref{w11:sec:reuse}).
  \item Sample splitting removes this bias by separating nuisance
    estimation from score evaluation, but wastes data
    (Section~\ref{w11:sec:splitting}).
  \item Cross-fitting recovers efficiency by rotating training
    and evaluation roles across $K$ folds, so each observation is
    evaluated only with nuisance estimates trained on other
    observations (Section~\ref{w11:sec:crossfit}).
  \item The cross-fitted AIPW estimator $\hat\tau_{\mathrm{DML}}$
    is the standard double machine learning estimator for the ATE
    (Section~\ref{w11:sec:dml}).
  \item Under the product-rate condition
    \eqref{w11:eq:product-rate}, orthogonality and cross-fitting
    together yield asymptotic linearity, root-$n$ consistency,
    asymptotic normality, and efficiency
    (Theorem~\ref{w11:thm:asymp}).  A sufficient symmetric case
    is that each nuisance estimator converges at least at rate
    $o_p(n^{-1/4})$ in mean square.
  \item The asymptotic variance is estimated consistently from the
    empirical variance of the out-of-fold estimated influence
    values (Theorem~\ref{w11:thm:var}).
  \item Machine learning improves nuisance estimation, but it does
    not resolve identification problems: unmeasured confounding,
    positivity failures, and untestable structural assumptions
    remain the researcher's responsibility
    (Section~\ref{w11:sec:limits}).
\end{enumerate}

% ------------------------------------------------------------
%  Key objects
% ------------------------------------------------------------
\begin{tcolorbox}[defstyle, title={Key Objects in This Chapter}]
\begin{center}
\renewcommand{\arraystretch}{1.4}
\begin{tabular}{lp{8.5cm}}
\hline
\textbf{Symbol} & \textbf{Meaning} \\
\hline
$\tau$
  & ATE $= \E\{Y(1)-Y(0)\}$ \\
$\mu_t(x)$
  & Outcome regression $\E(Y\mid T=t,X=x)$ \\
$\pi(x)$
  & Propensity score $P(T=1\mid X=x)$ \\
$\eta$
  & Nuisance tuple $(\mu_0,\mu_1,\pi)$ \\
$\phi(O;\,\psi,\eta)$
  & Generic orthogonal score; zero mean at truth \\
$\varphi_{\mathrm{eff}}(O;\,\tau,\eta)$
  & Efficient influence function for the ATE \eqref{w11:eq:orth-score-ate} \\
$\mathcal{I}_k$
  & $k$-th fold; $k=1,\ldots,K$ \\
$\hat\eta^{(-k)}$
  & Nuisance estimates trained without fold $k$ \\
$\hat\tau_{\mathrm{DML}}$
  & Cross-fitted AIPW (DML) estimator \eqref{w11:eq:dml-ate} \\
$\hat\varphi_i$
  & Out-of-fold estimated influence value for observation $i$ \\
$\hat{V}$
  & Variance estimator \eqref{w11:eq:var-est} \\
\hline
\end{tabular}
\end{center}
\end{tcolorbox}


% ============================================================
\section*{Exercises}
\label{w11:sec:exercises}
% ============================================================

\begin{enumerate}

\item \textbf{Verifying orthogonality.}
  \begin{enumerate}
    \item Write out the efficient influence function
      $\varphi_{\mathrm{eff}}(O;\,\tau,\eta)$ for the ATE and
      compute its expectation.  Confirm it equals zero at the truth.
    \item Consider a perturbation $\pi \mapsto \pi + r\cdot h$ for a
      bounded function $h(x)$.  Differentiate
      $\E\{\varphi_{\mathrm{eff}}(O;\,\tau, \mu_0, \mu_1, \pi+rh)\}$
      with respect to $r$ and evaluate at $r=0$.  Show the
      derivative is zero, confirming Neyman orthogonality in the
      propensity score direction.
    \item Repeat the calculation for a perturbation in $\mu_1$.
  \end{enumerate}

\item \textbf{First-order bias of the plug-in estimator.}
  Suppose $\hat\mu_1 = \mu_1 + \delta$ for a deterministic function
  $\delta(x)$, and $\hat\mu_0 = \mu_0$, $\hat\pi = \pi$.
  \begin{enumerate}
    \item Show that the bias of the resulting prediction estimator
      $\hat\tau_{\mathrm{pred}} = n^{-1}\sum_i
      \{\hat\mu_1(X_i) - \hat\mu_0(X_i)\}$ is
      $\E\{\delta(X)\}$.
    \item Show that the population bias of the plug-in AIPW
      estimator (using $\hat\mu_1 = \mu_1 + \delta$ with the
      true $\pi$ and $\mu_0$) is zero.  Explain which property
      of the AIPW score is responsible.  Does this mean the
      plug-in AIPW estimator is asymptotically unbiased when
      $\hat\mu_1$ is estimated flexibly from the same sample
      used to evaluate the score?  Explain why or why not.
  \end{enumerate}

\item \textbf{Cross-fitting with $K=2$.}
  Partition a sample of $n=200$ observations into two halves
  $\mathcal{I}_1$ and $\mathcal{I}_2$.
  \begin{enumerate}
    \item Write down the cross-fitted moment equation
      \eqref{w11:eq:crossfit-moment} explicitly for $K=2$.
    \item Explain why the contribution from $\mathcal{I}_1$
      uses nuisance estimates trained on $\mathcal{I}_2$ and
      vice versa.
    \item Compare this procedure to single sample splitting
      with $\mathcal{I}_1$ as the training fold.  What is
      gained and what is lost?
  \end{enumerate}

\item \textbf{The product-rate condition.}
  Suppose $\|\hat\pi - \pi\|_{L_2} = o_p(n^{-\alpha})$ and
  $\|\hat\mu_t - \mu_t\|_{L_2} = o_p(n^{-\beta})$ for
  $\alpha, \beta > 0$.
  \begin{enumerate}
    \item State the condition on $\alpha$ and $\beta$ under which
      the product-rate condition \eqref{w11:eq:product-rate} holds.
    \item Show that $\alpha = \beta = 1/4$ satisfies your condition.
    \item Does $\alpha = 1/2$, $\beta = 0$ satisfy it?  What does
      this say about the case where the propensity score is
      estimated parametrically but the outcome model is not
      consistent?
  \end{enumerate}

\item \textbf{Variance estimation.}
  Using the explicit formula \eqref{w11:eq:dml-ate}, write out
  the estimated influence value $\hat\varphi_i$ for a generic
  observation $i \in \mathcal{I}_k$.  Show that
  $n^{-1}\sum_i \hat\varphi_i = 0$ when $\hat\tau_{\mathrm{DML}}$
  solves the cross-fitted moment equation, and explain why the
  centering in \eqref{w11:eq:var-est} is therefore asymptotically
  inconsequential.

\item \textbf{What machine learning cannot do.}
  For each of the following scenarios, identify the identification
  assumption that is violated and explain why flexible nuisance
  estimation cannot remedy the problem.
  \begin{enumerate}
    \item A study of job-training effects on earnings omits
      pre-program earnings, a strong predictor of both
      program participation and subsequent earnings.
    \item A study of a medical treatment estimates the propensity
      score accurately, but the treatment was assigned in clusters
      (hospital wards) and outcomes may be correlated within
      clusters in a way that depends on the fraction of patients
      treated.
    \item An instrument is used to estimate the effect of
      education on wages, but the instrument also directly
      affects wages through a channel not related to education.
  \end{enumerate}

\end{enumerate}

% Bibliography (standalone only)
\ifSubfilesClassLoaded{%
  \bibliographystyle{plainnat}
  \bibliography{causal_inference}
}{}

\end{document}